% ============================================================
%  Kubernetes Notes
%  Compiled with pdfLaTeX (TeXShop on macOS)
% ============================================================
\documentclass[11pt, a4paper]{article}

% --- Packages ---
\usepackage[T1]{fontenc}
\usepackage[utf8]{inputenc}
\usepackage{lmodern}           % cleaner fonts
\usepackage[margin=2.5cm]{geometry}
\usepackage{parskip}           % space between paragraphs instead of indent
\usepackage{hyperref}          % clickable links & TOC
\usepackage{xcolor}            % colour support
\usepackage{listings}          % code / YAML blocks
\usepackage{booktabs}          % nicer tables
\usepackage{enumitem}          % control over lists

% --- hyperref setup ---
\hypersetup{
    colorlinks = true,
    linkcolor  = blue!60!black,
    urlcolor   = blue!60!black,
    citecolor  = green!50!black,
    pdftitle   = {Kubernetes Notes},
    pdfauthor  = {Eric R. Fulmer},
}

% --- listings setup (for YAML / shell snippets) ---
\lstdefinestyle{shell}{
    backgroundcolor = \color{gray!10},
    basicstyle      = \ttfamily\small,
    breaklines      = true,
    frame           = single,
    framerule       = 0.4pt,
    rulecolor       = \color{gray!50},
    numbers         = none,
    xleftmargin     = 6pt,
    xrightmargin    = 6pt,
}
\lstset{style=shell}

% ============================================================
\begin{document}

% --- Title block ---
\begin{center}
    {\LARGE \textbf{Kubernetes Notes}} \\[0.5em]
    {\large Personal Reference} \\[0.3em]
    {\small Last updated: \today}
\end{center}

\tableofcontents
\bigskip
\hrule
\bigskip

% ============================================================
\section{Introduction}

The purpose of this document is, obviously, my Kubernetes notes for myself. I am trying to avoid buzzwords and jargon, and explain both the \textit{hows} and \textit{whys} of Kubernetes. 

That implies two basic questions worth answering before we go any further:

\begin{itemize}
\item What \textit{is} Kubernetes?
\item Why does Kubernetes exist -- and why have so many people and organizations found it useful?
\end{itemize}

You'll often see Kubernetes described as "the operating system of the cloud." This is a buzzy way of describing it, but it is an accurate description. The different components of a Kubernetes cluster facilitate deploying applications in a cloud environment. That cloud may be public, like Amazon Web Services, or private, like your employer or school's server racks. 

Software suited for running on a cloud should be durable and largely agnostic to what computer it runs on, aside from fundamental incompatibilities on the level of x86 vs. ARM processor architectures, Linux vs. Windows vs. macOS, and so on. As operating systems provide convenient interfaces to the physical computer programs run on, they (ideally!) remove a programmer's need to know the specific CPU, RAM, and networking details of the machine their code runs on. Kubernetes does the same one level up -- abstracting away individual servers, nodes, regions, and storage types, handling concerns like application self-healing, upgrades and rollbacks, and storage on the programmer's behalf. This may not be a perfect analogy, but none is. It is the springboard for understanding those \textit{hows} and \textit{whys}.

Kubernetes was originally developed at Google, to make it easier for programmers to deploy and maintain programs at so-called "Google scale;" highly available, highly fault-tolerant systems. It is now owned and maintained by the Cloud Native Computing Foundation (CNCF). Before Kubernetes' release, containerization tools -- Docker being the most known -- were gaining notice and popularity for deploying industry applications. Docker handles a lot of the repeated complexities of deploying applications; packaging up everything you need to run the application, with a smaller footprint than a VM. However, Docker and Podman are \textit{containerization} tools. Their purpose is to set up a lightweight-but-isolated environment to run an application in. Things like application auto-healing and rollouts are not in scope for those programs. While Docker provides related tools for some of these, specifically Docker Compose and Docker Swarm, a solution more suitable for production deployments was still needed. Enter Kubernetes. 

% ============================================================
\section{Components of a Kubernetes Cluster}

A Kubernetes cluster consists of a \textbf{control plane node} (the "leader," "brain," or "main office" of the cluster) and one or more \textbf{worker nodes}. The control plane manages the desired state of the cluster; worker nodes run the actual workloads. Certain deployments, typically home lab or other sorts of non-production Kubernetes clusters, may run a single node which hosts both the control plane's components and runs workloads, but this is not advisable for production.

\subsection{API Server}

Coming soon!

\subsection{\texttt{etcd}}

A core component of Kubernetes is the durable \texttt{etcd} database. \texttt{etcd} is a key-value database (as opposed to a relational/SQL database, or a document store like MongoDB) used by many other distributed systems. \texttt{etcd} is a Kubernetes cluster's "source of truth" for the \textit{desired cluster state}; a concept that is crucial to understanding and using Kubernetes. 

Like most well-engineered pieces of software, access to the database is not a free-for-all affair. The Kubernetes API server is the intermediary between the \texttt{etcd} database and the outside world. The rest of this subsection discusses \texttt{etcd} and how Kubernetes uses it in greater depth; readers for whom this is not interesting or relevant can skip it. 

High-availability Kubernetes servers may want to run multiple \texttt{etcd} instances; I will go into more detail on this later. Clusters using multiple \texttt{etcd}s use the \textit{raft consensus algorithm} to designate a "leader" or main instance. 

\texttt{etcd} consists of three executables:

\begin{enumerate}
\item \texttt{etcd} - the program running the database itself.
\item \texttt{etcdctl} - provides a command-line interface for querying the database and handling admin tasks
\item \texttt{etcdutl} - used for tasks that require operating directly on \texttt{etcd} data files, like migrating data between versions of \texttt{etcd}, restoring snapshots, defragmenting the database, and database validation
\end{enumerate}

\subsubsection{How Kubernetes stores data in \texttt{etcd}}

\texttt{etcd} is a key-value database. Possible values for keys are seemingly pretty free-form, but Kubernetes uses a Linux path-like format for them. All keys start with \texttt{/registry/}, followed by the resource type, the namespace (if the object is namespaced), then the resource's name.

For example, the default namespace's definition is at \texttt{/registry/namespaces/default/}, and namespaced objects like Pods will live at, e.g. \texttt{/registry/pods/myAppNS/apiPod/}.

\subsubsection{Using \texttt{etcdctl}}

In most setups, \texttt{etcdctl} will require you to specify the API version, the endpoint and port the database server is listening at, and the certificate info. For a concrete example, if you're running a cluster locally using Docker Desktop, from within the Pod running \texttt{etcd}, your \texttt{etcdctl} commands may look something like

\begin{lstlisting}
ETCDCTL_API=3 etcdctl get /registry/namespaces/default \
  --endpoints=https://127.0.0.1:2379 \
  --cacert=/etc/kubernetes/pki/etcd/ca.crt \
  --cert=/etc/kubernetes/pki/etcd/server.crt \
  --key=/etc/kubernetes/pki/etcd/server.key
\end{lstlisting}

The values for the Kubernetes objects are stored in Protobuf format, so while this command should run successfully (at least on a cluster created by Docker Desktop), the output will not be human-readable. There are multiple ways to turn the Protobuf data into more readable formats, but the primary one is, of course, the Kubernetes API. 

To view the values in a human-readable form without leaving \texttt{etcdctl}, pass the \texttt{-w} flag (for "write," as in, write out to your terminal/display). Options include \texttt{fields}, \texttt{simple}, and \texttt{json}. \texttt{fields} and \texttt{json} will render the values for individual keys in a human-readable format, in my experience.

\subsubsection{Commands and Flags}
Here are some \texttt{etcdctl} commands and flags worth knowing:

\begin{itemize}
\item \texttt{etcdctl get|put|del} - see, set, or delete entries.
\item \texttt{etcdctl get {-}{-}prefix} - get all items whose key starts with the given prefix.
\item \texttt{etcdctl get {-}{-}keys-only} - get only the key, not the value, if it exists. Useful in conjunction with \texttt{{-}{-}prefix} to get every object of a specific type, in a specific namespace, or both.
\end{itemize}

\paragraph{Further reading}

\begin{itemize}
\item https://etcd.io/docs/v3.6/
\end{itemize}

\subsection{Controller Manager}

The \textit{controller manager} is a daemon that manages controllers. This implies a very obvious question -- what is a controller? Controllers programs that monitor the cluster (through the API server, Kubelet, or both?) to it to reconcile the desired and actual state of your cluster. Each controller manages a type or type of objects 

Additional controllers for Custom Resource Definitions (CRDs) can be added. 

\subsection{Scheduler}

Coming soon!

\subsection{Container Runtime}

Every Node, worker or Control Plane, must have at least one \textit{container runtime} installed on it. Container runtimes, such as Docker Engine, handle the crucial job of actually running containers. While Docker Engine is the most well-known, it can be heavyweight for some deployments. Other runtimes include containerd, CRI-O, and Mirantis Container Runtime, which is a commercially available runtime.

All container runtimes should be compatible with the Container Runtime Interface (CRI), which is an interface or protocol for a uniform method of communication between the kubelet and container runtime. 

CLI applications to interact with runtimes include \texttt{ctr}, \texttt{nerdctl}, and \texttt{crictl}. \texttt{crictl} is used for debugging purposes, \texttt{nerdctl} is a general-purpose Docker-like CLI application specifically for containerd, and \texttt{crictl} is another debugging-focused tool for all CRI-compatible runtimes, associated with Kubernetes.

Each of these applications connects to Unix sockets suck as:

\begin{itemize}
\item \texttt{unix:///var/run/dockershim.sock}
\item \texttt{unix:///var/run/containerd/containerd.sock}
\item \texttt{unix:///var/run/crio/crio.sock}
\item \texttt{unix:///var/run/cri-dockerd.sock}
\end{itemize}

\subsection{Kubelet}

The kubelet is an application that runs on each Node, with three major responsibilities:

\begin{enumerate}
\item registering the Node with the Control Plane upon Node start-up
\item managing communication between the Node and the Control Plane
\item managing Pods so that the container runtime can run their constituent containers, by doing things like:
	\begin{enumerate}
	\item mounting Volumes, ConfigMaps, Secrets
	\item running liveness and readiness probes, to monitor Pod health
	\item checking they respect any resource limits defined for them
	\end{enumerate}
\end{enumerate}

The kubelet communicates with the Control Plane primarily through the API Server, watching for any updates to PodSpecs (the specification component of a Pod manifest) scheduled to its specific Node. It also posts updates back to the API Server about the \textit{actual state} of those \textit{desired resources}.

The kubelet also runs its own small HTTPS API, through which it proxies requests that operate on the container-level, like \texttt{kubectl logs} and \texttt{kubectl exec}.

Aside from changes to the cluster state requested through \texttt{kubectl}, which are assigned to a Node by the scheduler and communicated to the kubelet via the API Server, the kubelet can also manage \textit{static Pods}. More on those later.

The kubelet is typically run as a daemon with \texttt{systemctl}, since it cannot manage itself as a Pod.

\subsection{Kube-Proxy}

Coming soon!

\end{document}